%% ============================================================
%%  arXiv submission source.  Compile with pdflatex (twice).
%%  No placeholders remain.  Worth a glance before posting: the MSC codes in
%%  \subjclass, and the arXiv categories (math.AG primary, cs.CC cross-list).
%% ============================================================
\documentclass[11pt,reqno]{amsart}
\usepackage{amsmath,amssymb,amsthm,mathtools}
\usepackage[margin=1.15in]{geometry}
\usepackage[colorlinks=true,linkcolor=black,citecolor=black,urlcolor=blue]{hyperref}

\theoremstyle{plain}
\newtheorem{theorem}{Theorem}[section]
\newtheorem{proposition}[theorem]{Proposition}
\newtheorem{lemma}[theorem]{Lemma}
\newtheorem{corollary}[theorem]{Corollary}
\theoremstyle{definition}
\newtheorem{definition}[theorem]{Definition}
\newtheorem{remark}[theorem]{Remark}
\newtheorem{question}[theorem]{Question}
\newtheorem{conjecture}[theorem]{Conjecture}

\DeclareMathOperator{\GL}{GL}
\DeclareMathOperator{\SL}{SL}
\DeclareMathOperator{\Sym}{Sym}
\DeclareMathOperator{\mult}{mult}
\DeclareMathOperator{\dfc}{def}
\DeclareMathOperator{\tr}{tr}
\DeclareMathOperator{\per}{per}
\DeclareMathOperator{\Ad}{Ad}
\newcommand{\cO}{\overline{\Omega}}
\newcommand{\bC}{\mathbb{C}}
\newcommand{\bN}{\mathbb{N}}
\newcommand{\bZ}{\mathbb{Z}}

\title[Conductors of orbit closures and the invariant of $\det_3$]
      {Conductors of orbit closures, and the fundamental invariant
       of the $3\times 3$ determinant}

\author{Swami Sethuraman}
\address{Beneficus AI}
\email{swsethuraman@beneficus.ai}

\date{\today}

\subjclass[2020]{Primary 14L24; Secondary 13A50, 20G05, 68Q17}
\keywords{orbit closure, geometric complexity theory, fundamental invariant,
determinantal complexity, conductor, boundary deficit}

\begin{document}

\begin{abstract}
Let $v$ be a form characterised by its symmetries, with finite stabiliser $H$
and affine orbit closure $\cO$.  The \emph{boundary deficit}
$\dfc(\lambda,\delta)=\dim S_\lambda^{H}-\mult_\lambda\bC[\cO]_\delta\ge 0$
measures the gap between the stabiliser's Peter--Weyl count and the coordinate
ring of the closure, and the \emph{conductor} $c(\lambda,\delta)$ is the depth
of that gap along the boundary, measured in multiplications by a boundary
semiinvariant.  We compute this invariant exactly in the two smallest
power-sum worlds: a closed form in the first, and in the second a transport
formula computing the conductor from boundary data alone, as the maximal
boundary-torus weight divided by the normal weight --- a shape we conjecture
holds generally on weights of non-empty support, and which we show can fail
without that hypothesis.  We then compute it for the $3\times 3$ determinant.

Our main results concern $\det_3$.  We prove that the minimal degree of a
nonzero $\SL_9$-invariant on $\cO=\overline{\GL_9\cdot\det_3}$ is
$e(\det_3)=18=2n^{2}$, with a one-dimensional space of such invariants, and we
evaluate the invariant exactly: it is nonzero at $\det_3$, with value
$-2^{16}3^{7}5^{3}7^{2}$ in a fixed integral normalisation.  B\"urgisser and Ikenmeyer
prove $e(\det_n)\ge n^{2}$ for every $n$, strictly when $n$ is odd, and their
degree-period theorem confines $e(\det_3)$ to the multiples of $6$ that are at
least $12$; equality $e(\det_n)=n^{2}$ is known only at $n=2$ and $n=4$, in
each case by an Alon--Tarsi-type nonvanishing that has no odd analogue.  Ours
is therefore the third determinant for which the fundamental invariant is
determined, and the first whose degree exceeds $n^{2}$; in particular
$e(\det_3)=18>16=e(\det_4)$, so the invariant degree is not monotone in $n$.  The lower bound
$e(\det_3)\ge18$ turns out to need no computation: an elementary parity count
on bracket monomials, which recovers Howe's vanishing theorems and reproduces
the classical first-invariant degrees for ternary cubics and cubic surfaces,
already excludes every smaller degree.  We show $V(\Phi_{18})\cap\cO=\partial\cO$ exactly, so that
degeneracy of a cubic of border determinantal complexity at most $3$ is
precisely the vanishing of one explicit degree-$18$ equation; we deduce that
the invariant ring of $\cO$ is a numerical semigroup ring, and, as a
demonstration of the method rather than a new result, we recover
$\per_3\notin\cO$ by evaluating a single integer.  We then exhibit the second
generator of that semigroup ring: the degree period forces an invariant in a
degree divisible by $6$ but not by $18$, the ambient space of degree-$24$
invariants is one-dimensional --- which makes vanishing as cheap to decide as
nonvanishing --- and
$\Phi_{24}(\det_3)=-2^{12}3^{7}5^{3}7^{2}\cdot11\cdot41\neq0$.  Hence
$\mathcal E(\det_3)\supseteq\langle 18,24\rangle$, its complement in $6\bN$ is
contained in $\{6,12,30\}$, and degree $30$ is the only value left
undetermined.  Finally we compute the
first conductor of a determinant: $c\big((2,2,2),2\big)=1$, certified at two
provably inequivalent orbit points and complete along the entire
$\Phi_{18}$-ray.  The conductor computation rests on an evaluation functional
several of whose natural symmetries we rule out, but whose totals we then
determine exactly: they are governed by a gauge $\Psi$, the restriction of the
fundamental degree-$6$ invariant of $\bC^{3}\otimes\bC^{3}\otimes\bC^{3}$ to
slab normal form.  The parabolic character controlling the functional is
$\det^{2}$ in the net's parameter slot, for the arithmetic reason that the two
parts of the relevant weight differ by $2$, and a one-dimensional uniqueness
statement then forces $\mathrm{TOTAL}=\Psi\cdot 1\,152\,144\,000$.
\end{abstract}

\maketitle

\section{Introduction and setup}\label{sec:setup}

Let $G=\GL(V)$ act on $W=\Sym^{d}V^{*}$, let $v\in W$ have finite stabiliser
$H\subset G$, and let $\cO=\overline{G\cdot v}$ be the affine orbit-closure
cone.  Algebraic Peter--Weyl gives the multiplicities on the orbit,
\[
  \mult_\lambda\bC[G\cdot v]=\dim (S_\lambda^{*})^{H}=:m(\lambda),
\]
a purely group-theoretic quantity.  The multiplicities of the \emph{closure}
are the hard side, and they are the side that carries the obstruction theory
of geometric complexity theory \cite{MS,BIP,IK}.  Landsberg's book
\cite{Lands} is the standard account of the geometry involved, and the
programme it describes --- separating $\overline{\GL\cdot\per_m}$ from
$\overline{\GL\cdot\det_n}$ by representation-theoretic data --- is what makes
the closure multiplicities worth the trouble; the occurrence obstructions
originally proposed for the task are now known not to suffice \cite{BIP}, which
sharpens rather than removes the interest in computing the closure side
exactly.

Suppose the boundary $\cO\setminus G\cdot v$ is cut set-theoretically by a
semiinvariant $\Delta$ of degree $e$ and weight $\det^{w}$.  Multiplication by
$\Delta$ embeds isotypic pieces, so along the ray
$(\lambda+mw\mathbf 1,\ \delta+me)$ the closure multiplicity is nondecreasing
and stabilises at $m(\lambda)$, since $\bC[G\cdot v]=\bC[\cO][1/\Delta]$.

That the deficit is generally nonzero is not a technicality but the central
difficulty of the subject: the orbit closures of $\det_n$ and $\per_n$ fail to
be normal for every $n>2$ \cite{KumarCMH}, quantitative control on invariants
of such closures is delicate \cite{DM}, and it is precisely this failure
that obstructs the Mulmuley--Sohoni programme's passage from the orbit, where
multiplicities are group theory, to the closure, where they are not.  The
deficit and the conductor are an attempt to measure that failure rather than
work around it.

\begin{definition}\label{def:deficit}
The \emph{boundary deficit} is
$\dfc(\lambda,\delta)=m(\lambda)-\mult_\lambda\bC[\cO]_\delta\ \ge 0$,
and the \emph{conductor} $c(\lambda,\delta)$ is the least $k$ with
$\dfc(\lambda+kw\mathbf 1,\delta+ke)=0$: the maximal pole order along the
boundary among functions in the $\lambda$-isotypic component.
\end{definition}

\begin{remark}[On the word]\label{rem:name}
The classical conductor is the ideal $\mathfrak c=\operatorname{Ann}_{O}(\tilde
O/O)$ of a ring in its normalisation, and non-normality of orbit closures is
active in this very context, so the choice of name needs justifying.  We do not
mean $\mathfrak c$: here $\bC[G\cdot v]$ is a \emph{localisation} of
$\bC[\cO]$, not its normalisation, and being non-finite over it, the
corresponding conductor ideal is zero.  What Definition \ref{def:deficit}
measures is the conductor in the numerical-semigroup sense --- the least index
from which membership holds and continues to hold --- and this is literal
rather than analogical: Corollary \ref{cor:semigroup} shows the invariant ring
of $\cO$ \emph{is} a numerical semigroup ring, and along each ray the deficit
vanishes from $c(\lambda,\delta)$ onwards.  Equivalently, $c$ is the least $k$
with $\Delta^{k}\,\bC[G\cdot v]_\lambda\subseteq\bC[\cO]$, which is the
conductor of $\bC[\cO]$ in the module $\bC[G\cdot v]$ restricted to one
isotypic piece.
\end{remark}

That such a $k$ exists, at an explicit uniform bound, is a theorem of
Ikenmeyer and Kandasamy \cite[Lem.~5.2]{IK} for power sums in the regime
(number of variables) $\ge$ (degree); their argument is algebraic and makes no
optimality claim.  Everything in this paper concerns the exact minimal index
and the boundary geometry that computes it.

Sections \ref{sec:worldA} and \ref{sec:worldB} determine the invariant in the
two smallest classical cases, where every ingredient is available in closed
form; these serve both as results and as calibration for the determinant.
Section \ref{sec:det} contains the main results.  Section
\ref{sec:functional} analyses the evaluation functional used there, including
what we could and could not establish about it.  Section \ref{sec:method}
describes the verification protocol; we consider this part of the content
rather than an appendix, since several of the results are computational.

\section{Binary quartics: a closed form}\label{sec:worldA}

Take $W=\Sym^{4}\bC^{2}$, $v=x^{4}+y^{4}$, $H=\mu_4^{2}\rtimes S_2$ of order
$32$, so that $\cO$ is the catalecticant cubic hypersurface $\{J=0\}$.  Since
$\operatorname{disc}=256(I^{3}-27J^{2})$ identically, $\operatorname{disc}$
restricted to $\cO$ equals $256\,I^{3}$: the boundary divisor is cut by the
quadratic invariant $I$ (degree $2$, weight $\det^{4}$), along which the
discriminant ramifies with index $3$, and it is $I$, not the discriminant, that
generates the correct ray.  The dense boundary orbit $G\cdot x^{3}y$ has
stabiliser torus $\operatorname{diag}(t,t^{-3})$; the boundary point is a
smooth point of $\cO$ with normal weight of size $8$.

\begin{theorem}\label{thm:A}
For every weight $(a,b)$ with $a+b=4\delta$,
\[
  \dfc\big((a,b),\delta\big)=c\big((a,b),\delta\big)
   =\max\Big(0,\Big\lfloor\tfrac{a-3b}{8}\Big\rfloor\Big),
\]
equivalently
$\mult_{(a,b)}\bC[\cO]_\delta=N(a,b)-\lfloor (a-3b)/8\rfloor^{+}$, where
$N(a,b)$ counts unordered pairs $\{p,q\}$ with $p+q=a-b$ and
$p\equiv q\equiv -b \pmod 4$, excluding $p=q$ when $b$ is odd.  Consequently
$\dfc>0$ if and only if $b\le\delta-2$; the maximal conductor at degree
$\delta$ is $\lfloor\delta/2\rfloor$; and the total deficit at degree $\delta$
is $\lfloor\delta^{2}/4\rfloor$.
\end{theorem}

The proof has two halves.  The pole law is computed in the Sylvester model on
the transversal family $f_s=x^{3}y+s\,xy^{3}$, which lies in $\{J=0\}$
identically with $I(f_s)=-s/4$; the Waring pair is $c=\pm 1/(8\sqrt s)$,
$\ell=x\pm\sqrt s\,y$, and the isotypic component is spanned by the classes
$e_{p,q,b}=\Sym[\hat\ell_1^{p}\hat\ell_2^{q}]\cdot(\hat\ell_1\wedge\hat\ell_2)^{b}$.
Writing the bracket coefficient at $u^{k}$ as
$A_k(p,q)\big[\iota^{q}+(-1)^{b+k}\iota^{p}\big]$ with
$A_k=[u^{k}](1+u)^{p}(1-u)^{q}$ and $\iota^{4}=-1$, the phase survives exactly
when $k\equiv b+(p-q)/4\pmod 2$, which coincides with the integrality parity;
so only $k_0\in\{0,1\}$ occurs, and $A_0=1$, $A_1=p-q\ne 0$ away from the
identically vanishing class.  Hence
$\operatorname{pole}(e_{p,q,b})=\lfloor (p+q-2b)/8\rfloor=\lfloor(a-3b)/8\rfloor$,
uniformly over the classes at a given weight.  Second, $\GL_2$-torus weight
spaces are one-dimensional, so all leading terms at a given pole level are
proportional to a single monomial and the pole filtration climbs by at most one
per level; combined with the identity above, deficit equals conductor equals
the pole value.  The closed forms follow, the total by
$\sum_{b\le\delta-2}\lfloor(\delta-b)/2\rfloor=\lfloor\delta^{2}/4\rfloor$;
this last identity is independently established as an identity of rational
functions.

\begin{remark}
The stratified picture closes completely here.  The tangent developable's
normalisation is $\bigoplus_\delta \Sym^{3\delta}\otimes\Sym^{\delta}$ under
$(\ell,m)\mapsto \ell^{3}m$ \cite{AFPRW}, one copy of each $(4\delta-b,b)$ for
$b\le\delta$; the non-normality gap is exactly one copy of
$S_{(4\delta-1,1)}$ per degree --- the equivariant globalisation of
$k[t^{2},t^{3}]\subset k[t]$ along the cuspidal edge.  The full tower reads
$\lfloor(a-3b)/8\rfloor^{+}\to[b=1]\to 0$.
\end{remark}

\section{Ternary cubics: a transport formula}\label{sec:worldB}

Take $W=\Sym^{3}\bC^{3}$, $v$ the Fermat cubic, $H=\mu_3^{3}\rtimes S_3$ of
order $162$, so that $\cO$ is the Aronhold quartic hypersurface $\{S=0\}$.  The
space of degree-$12$ invariants on $\cO$ is one-dimensional, whence
$\operatorname{disc}|_{\cO}=c\,T^{2}$ for the canonical degree-$6$
semiinvariant $T$.  The dense boundary orbit is the cuspidal locus
$G\cdot(x^{2}y+z^{3})$: a smooth point of $\cO$ with stabiliser torus
$\operatorname{diag}(t,t^{-2},1)$, normal weight $-6$, and
$\nu_{\partial}(T)=1$ exactly.

\begin{theorem}\label{thm:B}
On every weight with $m(\lambda)>0$ and $\delta\le 10$ --- all $254$ of them ---
\[
  c(\lambda)=\Big\lfloor\frac{\lambda_1-2\lambda_3}{6}\Big\rfloor
           =\Big\lfloor\frac{\mu_{\max}(\lambda)}{|w_N|}\Big\rfloor ,
\]
the maximal boundary-torus weight over the normal weight.  The inequality
$\le$ holds in general by the block-order estimate below; the equality is
verified computationally on the stated range.  The hypothesis is stated as
$m(\lambda)>0$ rather than as positivity of the deficit because $m(\lambda)$ is
a stabiliser character count --- cheap, and computable without any knowledge of
$\cO$ --- so the theorem can be applied where one does not already know the
answer.  Remark \ref{rem:orphan} shows the hypothesis cannot simply be dropped.
\end{theorem}

On the transversal family the normalised Waring blocks have exact orders
$-1$ for the two colliding lines, $0$ for the untouched line, $+1$ for the
collision wedge, $-1$ for mixed wedges and $+1$ for the determinant.  A class
with single-counts $(a_1,a_2,a_3)$, wedge-counts $(q_{12},q_{13},q_{23})$ and
determinant power $r=\lambda_3$ therefore has pole
$(a_1+a_2+q_{13}+q_{23}-q_{12}-r)/6$, whose maximum over class types is
$(\lambda_1-2\lambda_3)/6$.  Since Young symmetrisation only cancels and the
first fundamental theorem gives spanning, the bound is rigorous.

\begin{remark}[Orphans]\label{rem:orphan}
Attainment fails exactly on weights of empty support, where the contraction
shadow reports a positive pole while the deficit is $0$: the Young projection
kills the top.  In the sweep over all $\lambda$ with $\delta\le10$ and
$\lambda_1\le20$ the weights with $m(\lambda)=0$ and shadow pole $\ge1$ are
\[
  (10,1,1)\ \ [\delta=4],\qquad (13,1,1)\ \ [\delta=5],\qquad
  (11,11,2)\ \ [\delta=8],
\]
and extending to $\delta\le16$ adds $(17,17,5)$ at $\delta=13$.  They are not
independent: for $\GL_3$,
\[
  V_{(10,1,1)}^{*}\otimes\det^{12}=V_{(11,11,2)},
  \qquad
  V_{(13,1,1)}^{*}\otimes\det^{18}=V_{(17,17,5)},
\]
with Weyl dimensions $55,55$ and $91,91$ respectively.  \emph{The orphans come
in dual pairs.}  We do not know whether the orphan locus is stable under
$\lambda\mapsto\lambda^{*}\otimes\det^{k}$ in general; if it is, it is a
statement about where attainment fails, and bears directly on Question
\ref{q:transport}.  Thus the formula computes the shadow maximum, and on empty
support the projection cancels it.
\end{remark}

\begin{remark}[Two kinds of tower]\label{rem:towers}
The increments along the $T$-ray are multiplicities of $\bC[\cO]/(T)$, the
coordinate ring of the cusp cone, bounded by the cusp stabiliser's invariant
counts, with the shortfall being the cusp cone's own deficit: the structure
recurses down the stratification.  In the binary-quartic case the recursion
terminates after one step.  In the ternary-cubic case it does not: the level-two
correction admits no floor law and passes to the conic-with-tangent stratum,
whose stabiliser is two-dimensional and non-reductive~\cite{Kraft}.  The
deficit there is a lattice count rather than a floor, while the conductor
remains a floor --- the sharpest structural difference between the two cases.
The first case closes level by level only because every parametrisation in its
tower is finite.
\end{remark}

\section{The determinant}\label{sec:det}

Let $n=3$, $V=\bC^{9}$ with coordinates the entries of a $3\times3$ matrix, and
$v=\det_3\in \Sym^{3}V^{*}$.  Write $G=\GL_9$, $\cO=\overline{G\cdot\det_3}$,
and let $H\subset G$ be the stabiliser: by Frobenius,
$H=\big((\GL_3\times\GL_3)/\bC^{*}\big)\rtimes\bZ/2$ acting by
$X\mapsto PXQ$ and $X\mapsto X^{t}$, of dimension $17$.  This $H$ is
\emph{not} finite, so the deficit of Definition \ref{def:deficit} is taken with
$m(\lambda)=\dim(S_\lambda^{*})^{H}$ as before; the Peter--Weyl side is now a
branching count rather than a character count, but the definition and the
inequality $\dfc\ge0$ are unchanged.

\subsection*{The fundamental invariant}
An $\SL_9$-invariant of degree $\delta$ on $\Sym^{3}\bC^{9}$ has $\GL_9$-weight
$\det^{\delta/3}$ and corresponds to the rectangular weight
$\lambda=\big((\delta/3)^{9}\big)$; in particular $3\mid\delta$.  Following
\cite{BI} we call
\[
  e(v)\;=\;\min\{\delta>0:\ \bC[\overline{G\cdot v}]^{\SL_9}_\delta\neq0\}
\]
the degree of the fundamental invariant of the orbit closure, and we write
$\mathcal E(v)$ for the degree monoid $\{\delta:\ \bC[\overline{G\cdot
v}]^{\SL_9}_\delta\neq0\}$.

Three published facts narrow the search before any computation.  Howe's theorem
\cite[Thm.~3.14]{BI} gives $\bC[\Sym^{3}\bC^{9}]^{\SL_9}_\delta=0$ for
$\delta<9$, and gives $0$ again at $\delta=9$ because the form degree $3$ is
odd; hence $e(\det_3)>9$ \cite[Thm.~3.15]{BI}.  Next, $\det_3$ is polystable
\cite[Cor.~2.9]{BI}, so the degree monoid generates $b\,\bZ$ for the degree
period $b=(m/D)\,a$ \cite[Thm.~3.4]{BI}, where $a$ is the stabiliser period;
since $3\equiv3\pmod 4$ we have $a(\det_3)=2$ \cite[Thm.~2.5]{BI} and
therefore
\[
  b(\det_3)=\tfrac{9}{3}\cdot 2=6,
  \qquad
  \mathcal E(\det_3)\subseteq 6\bZ,
  \qquad
  \gcd\mathcal E(\det_3)=6 .
\]
So the candidate degrees are $12,18,24,\dots$, and the weights $\delta=3,9,15$
carry no invariant for reasons of period alone.

A second, entirely elementary constraint removes the rest of the range.  It is
a parity count on bracket monomials, and it extends the argument by which Howe
disposes of $\delta=m$ when $D$ is odd; we state it in general because the
statement costs nothing extra and because it is what does the work here.

\begin{proposition}[Bracket census]\label{prop:census}
Let $D$ be odd and let $\delta,m$ satisfy $m\mid \delta D$.  Put $k=\delta D/m$.
If
\[
  \delta \;>\; \binom{k}{D},
\]
then $\bC[\Sym^{D}\bC^{m}]^{\SL_m}_{\delta}=0$; consequently
$\bC[\overline{G\cdot v}]^{\SL_m}_{\delta}=0$ for every $v$.
\end{proposition}

\begin{proof}
By the first fundamental theorem for $\SL_m$ \cite{Weyl}, the space of
invariants of degree $\delta$ is spanned by products of $m\times m$ brackets in
$\delta$ symbolic letters, each letter occurring exactly $D$ times.  A bracket
is a determinant, so its $m$ letters are distinct; counting incidences,
$\delta D=km$ and each letter occupies exactly $D$ of the $k$ brackets.  Record
this as a $0/1$ matrix with rows indexed by letters, row sums $D$ and column
sums $m$.

Suppose two letters $i\ne j$ have equal rows, so that they occur together in
the same $D$ brackets and in no others.  The transposition exchanging them
fixes every bracket of the monomial setwise and swaps two arguments in each of
the $D$ brackets containing both, multiplying the monomial by $(-1)^{D}=-1$.
But letters are interchangeable placeholders for one and the same form, so the
transposed monomial evaluates to the \emph{same} invariant.  That invariant is
therefore its own negative, hence zero.

The rows are $D$-subsets of $[k]$, of which there are $\binom{k}{D}$.  If
$\delta>\binom{k}{D}$ two rows coincide for every monomial, so every spanning
monomial evaluates to zero.  The statement for an orbit closure follows because
$\SL_m$ is reductive, so restriction
$\bC[\Sym^{D}\bC^{m}]^{\SL_m}_\delta\to\bC[\overline{G\cdot v}]^{\SL_m}_\delta$
is surjective.
\end{proof}

Two checks, both against classical answers.  For ternary cubics ($m=D=3$) the
bound admits $\delta=4$ and excludes $\delta\le3$: the first invariant is
Aronhold's $S$, in degree $4$.  For cubic surfaces ($m=4$, $D=3$) it admits
$\delta=8$ and excludes $\delta=4$: the invariant ring of the cubic surface is
generated in degrees $8,16,24,32,40,100$.  The bound also recovers Howe's two
statements --- for $\delta<m$ no bracket monomial exists at all, and for
$\delta=m$ with $D$ odd every letter lies in all $k=D$ brackets, so all rows
coincide.

For $\det_3$ we have $m=9$, $D=3$, $k=\delta/3$, and the condition
$\delta\le\binom{\delta/3}{3}$ first holds at $\delta=18$:
\[
\begin{array}{c|cccccc}
 \delta & 3 & 6 & 9 & 12 & 15 & 18\\ \hline
 k=\delta/3 & 1 & 2 & 3 & 4 & 5 & 6\\
 \binom{k}{3} & 0 & 0 & 1 & 4 & 10 & 20\\
 \delta\le\binom{k}{3} & \text{no} & \text{no} & \text{no} & \text{no} &
   \text{no} & \text{yes}
\end{array}
\]

\begin{corollary}\label{cor:lower}
$e(\det_3)\ge 18$, with no computation.
\end{corollary}

At $\delta=18$ the census is silent, and one exact count settles what it leaves
open.

\begin{theorem}[Ambient census]\label{thm:ambient}
$\dim_\bC\bC[\Sym^{3}\bC^{9}]^{\SL_9}_{18}=1$.
\end{theorem}

\begin{proof}
The dimension is the plethysm coefficient
$\langle h_{18}[h_3],\,s_{(6^{9})}\rangle$, evaluated by expanding
$h_{18}[h_3]$ in the power-sum basis and pairing against the character
$\chi^{(9^{6})}$ by the Murnaghan--Nakayama rule.  The expansion has $44\,237$
distinct classes and the value is $1$.  The computation is exact rational
arithmetic throughout and is independent of everything in Sections
\ref{sec:functional} and \ref{sec:method}.
\end{proof}

\begin{definition}\label{def:phi18}
$\Phi_{18}$ denotes a nonzero element of $\bC[\Sym^{3}\bC^{9}]^{\SL_9}_{18}$.
By Theorem \ref{thm:ambient} it is determined up to a scalar, and every
statement we make about it below --- non-vanishing, the divisor, the semigroup,
and all ratios of values --- is independent of that scalar.  Where a bare
numerical value appears, the normalisation in force is stated with it.
\end{definition}

Three objects in this circle of ideas carry the letter $\Phi$ or $\phi$, and it
is worth separating them.  B\"urgisser and Ikenmeyer write $\Phi_w$ for the
fundamental invariant of a polystable form $w$, of degree $e(w)$ and normalised
by $\Phi_w(w)=1$; their subscript names the \emph{form}, ours names the
\emph{degree}, and the two agree up to scale,
$\Phi_{18}=\Phi_{18}(\det_3)\cdot\Phi_{\det_3}$.  They also write $\phi_w$ for
the invariant of degree $b(w)$ defined on the \emph{orbit} alone by
$g\cdot w\mapsto\det(g)^{a(w)}$ \cite[Lem.~3.2]{BI}; here $b=6$, and since
$\bC[\cO]^{\SL_9}_{18}$ is one-dimensional,
\[
  \Phi_{18}\big|_{G\cdot\det_3}\;=\;c\,\phi^{3}
  \qquad\text{for some }c\in\bC^{\times}.
\]
The degree monoid of Corollary \ref{cor:semigroup} is exactly
$\{6k:\phi^{k}\text{ extends regularly to }\cO\}$, so a question about which
degrees carry invariants is a question about which powers of $\phi$ extend
across the boundary.

This is what makes $\Phi_{18}$ an object rather than an output.  It is defined
by a dimension count, not by the algorithm that evaluates it, and its
uniqueness has three consequences we use.  First, $e(\det_3)\le18$ needs only
that $\Phi_{18}$ does not vanish identically on $\cO$, and a single nonzero
value at a single point of the orbit certifies that.  Second, the restriction
map $\bC[\Sym^{3}\bC^{9}]^{\SL_9}_{18}\to\bC[\cO]^{\SL_9}_{18}$ is either zero
or injective, so the census reduces $e(\det_3)$ to \emph{one bit}, and one
evaluation decides it.  Third, since the ambient space is one-dimensional, the
value of $\Phi_{18}$ at any point of $\Sym^{3}\bC^{9}$ --- including points
outside $\cO$ --- is canonical up to the single global scale, which
Corollary \ref{cor:perm} needs.

\begin{theorem}\label{thm:census}
$\dim_\bC\bC[\cO]^{\SL_9}_{18}=1$, and therefore
\[
  e(\det_3)=18=2n^{2}.
\]
\end{theorem}

\begin{proof}
$\ge$ is Corollary \ref{cor:lower}.  For $\le$, Theorem \ref{thm:ambient} and
Definition \ref{def:phi18} make $\Phi_{18}$ the unique candidate up to scale,
and $\Phi_{18}(\det_3)\neq0$ by \eqref{eq:value}; since $\det_3\in\cO$, the
restriction of $\Phi_{18}$ to $\cO$ is not the zero function.  One-dimensionality
of the restricted space follows from that of the ambient space.
\end{proof}

We are explicit about the dependencies, because the reader who has looked ahead
to Section \ref{sec:functional} will want to know how much of it Theorem
\ref{thm:census} rests on.  The answer is: none.  The lower bound is
combinatorial; the ambient count is a plethysm; the upper bound is one integer.
The gauge law of Section \ref{sec:functional}, which is verified and not proved,
and the covariance hypotheses refuted there, enter nowhere in this section.
Section \ref{sec:functional} studies \emph{how} the evaluation behaves across
orbit points; whether the invariant exists does not depend on it.

An earlier form of this argument established the vanishing at
$\delta=3,6,9,12,15$ by direct multiplicity computation, and those runs were in
fact carried out; the bracket census makes all five superfluous.  We have kept
them, and report that they return zero at every one of the five, because a
computation that agrees with a proof it did not use is the cheapest calibration
available and this one exercises the whole pipeline at the exact weights that
matter.  The vanishing statements are
exact multiplicity computations in
$\bC[\cO]_\delta$ at the rectangular weights, run in the streamed level
algorithm of Section \ref{sec:method}; the nonvanishing at $\delta=18$ is
witnessed by the value \eqref{eq:value}, which is a single exact integer and
therefore settles the question on its own.

\begin{theorem}\label{thm:value}
$\Phi_{18}(\det_3)\neq0$.  In the integral normalisation $\mathcal N$ --- the
evaluation functional of Section \ref{sec:functional} taken with unit weight on
the standard slab basis ---
\begin{equation}\label{eq:value}
  \Phi_{18}(\det_3)\;=\;-877\,879\,296\,000\;=\;-2^{16}\,3^{7}\,5^{3}\,7^{2} .
\end{equation}
\end{theorem}

\begin{remark}[What depends on the scale, and what does not]\label{rem:norm}
Only the non-vanishing is used in the proof of Theorem \ref{thm:census}; the
integer itself is a statement about $\mathcal N$.  Under B\"urgisser and
Ikenmeyer's convention $\Phi_{\det_3}(\det_3)=1$ holds by fiat and carries no
information, so the scale-invariant content of the two values computed here is
their ratio
\[
  \frac{\Phi_{18}(\per_3)}{\Phi_{18}(\det_3)}\;=\;-\frac{272}{4725},
\]
which is independent of every choice.  The canonical integral choice would be
the primitive generator of the rank-one lattice of integral invariants in this
degree, unique up to sign --- the same convention under which
$\Psi=-I_6^{\mathrm{prim}}$ in Section \ref{sec:functional}.  We have not
determined whether $\mathcal N$ is primitive.  The content divides the gcd of
the two values, namely $2^{16}3^{4}\cdot5\cdot7$, which is a weak bound; if
$\mathcal N$ is a proper multiple of the primitive generator then
\eqref{eq:value} must be divided by that multiple and the factorisation
restated.  We flag this because the factorisation, rather than the
non-vanishing, is the part a reader is likely to quote.
\end{remark}

\subsection*{Boundary geometry}

Two facts about this boundary are already known, and we use both.
H\"uttenhain and Lairez \cite{HL}, answering a question of Landsberg,
determine the boundary of the orbit of $\det_3$ completely (see also
\cite{Hue}): it is \emph{pure}
of codimension one with exactly two irreducible components, the orbit closures of the determinant of the generic
traceless matrix and of the universal quadric
$x_4x_1^{2}+x_5x_2^{2}+x_6x_3^{2}+x_7x_1x_2+x_8x_2x_3+x_9x_1x_3$.  We write
$P_1$ and $P_2$ for these two components (in \cite{HL} the same symbols denote
the two forms).  Independently, B\"urgisser and Ikenmeyer show for any
polystable $w$ that the zero set of the fundamental invariant in the orbit
closure is exactly the boundary, and that $\bC[G\cdot w]$ is the localisation
of $\bC[\overline{G\cdot w}]$ at it \cite[Prop.~3.9]{BI}.  Applied to
$\Phi_{18}$ this gives $V(\Phi_{18})\cap\cO=\partial\cO$ and puts
$\Phi_{18}$ in the role of the semiinvariant $\Delta$ of Section
\ref{sec:setup}, of degree $e=18$ and weight $\det^{6}$.  We verified the
vanishing directly at both of H\"uttenhain and Lairez's representatives, which
is a check on our normalisation rather than a new result.

What is left to compute is how $\Phi_{18}$ vanishes, and that is not formal.

\begin{proposition}\label{prop:divisor}
As a divisor on $\cO$,
\[
  \operatorname{div}\big(\Phi_{18}|_{\cO}\big)\;=\;6P_1+9P_2 .
\]
\end{proposition}

The two multiplicities arise differently, and the asymmetry is the point.  Along
$P_1$ the orbit closure is smooth at a generic point, a transversal
one-parameter family meets the divisor transversally, and the order is read off
from the weight: an $\SL_9$-invariant of degree $18$ satisfies
$\Phi_{18}(g\cdot v)=\det(g)^{6}\Phi_{18}(v)$, giving ramification $6$.
Along $P_2$ that argument is unavailable --- the orbit closure of $\det_n$ is
not normal for $n>2$ \cite{KumarCMH}, and $P_2$ is where the non-normality
sits --- and the multiplicity $9$ is obtained instead by a computation, from
integrality of the torus weights on the component.  We record it as a computed
value rather than a proved one; deriving it directly from the explicit
degeneration to $P_2$ that H\"uttenhain and Lairez supply --- via Jacobi's
formula, with $A$ skew-symmetric and $S$ symmetric --- is the obvious next step
and we have not carried it out.  In particular the multiplicities are not both
determined by the character, and a reader who expects $6$ on both components
should note that this is exactly what fails.

\begin{corollary}\label{cor:semigroup}
Restriction to the dense orbit is injective on $\bC[\cO]^{\SL_9}$ and
$\dim\bC[\cO]^{\SL_9}_\delta\le1$ for every $\delta$, so
$\bC[\cO]^{\SL_9}$ is the semigroup ring of the numerical semigroup
$S=\mathcal E(\det_3)$.  Combining Theorem \ref{thm:census} with the degree
period, $S\subseteq6\bN$, $\min S=18$ and $\gcd S=6$; hence $S\neq18\bN$.  By
Theorem \ref{thm:deg24}, $24\in S$, and therefore
\[
  S\;\supseteq\;\langle 18,24\rangle\;=\;6\cdot\langle3,4\rangle
  \;=\;\{0,18,24,36,42,48,54,\dots\} .
\]
Equivalently $\tfrac16S$ contains every nonnegative integer except possibly
$1,2,5$, and $1,2\notin\tfrac16S$.  The only undetermined degree is $30$:
either $S=\langle18,24\rangle$, or $S=6\cdot\langle3,4,5\rangle
=\{0,18,24,30,36,\dots\}$.
\end{corollary}

In the first version of this paper the last line of Corollary
\ref{cor:semigroup} was a pure existence statement: were $S=18\bN$ its greatest
common divisor would be $18$, while the degree period forces $6$, so $\cO$ had
to carry an invariant in some degree that is a multiple of $6$ but not of $18$,
and $24$ was the smallest candidate.  That invariant now exists explicitly.

\begin{theorem}\label{thm:deg24}
$\dim_\bC\bC[\Sym^3\bC^9]^{\SL_9}_{24}=1$; write $\Phi_{24}$ for a generator.
Its value on the determinant is nonzero:
\begin{equation}\label{eq:value24}
  \Phi_{24}(\det_3)\;=\;-24\,745\,222\,656\,000
  \;=\;-2^{12}\,3^{7}\,5^{3}\,7^{2}\cdot11\cdot41\;\neq\;0 ,
\end{equation}
in the normalisation fixed by the explicit bracket monomial recorded in
\texttt{results/results\_deg24.md}.  Consequently
\[
  24\in\mathcal E(\det_3),\qquad \dfc\big((8^{9}),24\big)=0,\qquad
  \mult_{(8^{9})}\bC[\cO]_{24}=1 ,
\]
and $\Phi_{24}$ is the second generator of $\bC[\cO]^{\SL_9}$ --- a new
generator, not a product, since there is no invariant of degree $6$ on $\cO$ to
multiply $\Phi_{18}$ by.
\end{theorem}

The one-dimensionality in Theorem \ref{thm:deg24} is what makes the statement
cheap, and it is worth isolating why.  For a general degree the vanishing
direction is the expensive one: it needs the multiplicity of the rectangular
weight in $\bC[\cO]_\delta$, a computation strictly larger than the
$\delta=18$ census.  Here it is free.  Since
$\bC[\cO]^{\SL_9}_{24}$ is the image of $\bC[\Sym^3\bC^9]^{\SL_9}_{24}$ under
restriction --- surjective on invariants because $\SL_9$ is reductive --- and
that ambient space is a single line, $24\in\mathcal E(\det_3)$ holds
\emph{if and only if} $\Phi_{24}(\det_3)\neq0$.  One integer decides both
directions.  The ambient Hilbert function, computed exactly, is
\[
  \dim\bC[\Sym^3\bC^9]^{\SL_9}_\delta \;=\;
  0,0,0,0,0,\;\mathbf 1,\;0,\;\mathbf 1,\;1,\;4
  \qquad(\delta=3,6,\dots,30).
\]
Two features deserve comment.  The bracket census of Proposition
\ref{prop:census} permits $\delta=21$ --- $21\le\binom73=35$ --- yet the
multiplicity there is $0$, so the exact count is strictly sharper than the
pigeonhole bound.  And at $\delta=30$ the ambient space is four-dimensional, so
the same collapse does not occur: degree $30$ is not decided by a single number,
and Corollary \ref{cor:semigroup} leaves it open for that reason.

\begin{remark}[A normalisation-free form of \eqref{eq:value24}]\label{rem:norm24}
The scale of $\Phi_{24}$ is a convention, so \eqref{eq:value24} should be read
alongside a ratio that is not.  Restrict $\Phi_\delta$ to the six-dimensional
space $U$ of cubics $\sum_{\sigma\in S_3}u_\sigma
x_{1\sigma1}x_{2\sigma2}x_{3\sigma3}$, which contains both $\det_3$ and
$\per_3$.  The diagonal torus of $\GL_9$ preserves $U$ and the weight condition
$\sum_\sigma e_\sigma M_\sigma=8J$ has exactly the nine solutions
$e=(a,a,a,8-a,8-a,8-a)$, so
$\Phi_{24}|_U=\sum_{a=0}^{8}K_aP^{a}Q^{8-a}$ with $P=u_{\mathrm{id}}u_cu_{c^2}$
and $Q=u_{t_1}u_{t_2}u_{t_3}$; a row/column permutation of odd total sign
exchanges $P$ and $Q$ and acts trivially on $\det^{8}$, so $K_a=K_{8-a}$.  Then
$\Phi_{24}(\det_3)=\sum_a(-1)^aK_a$ and $K_8$ is the value of the run in which
only even permutations are allowed, and
\[
  \Phi_{24}(\det_3)/K_8\;=\;17325\;=\;3^{2}\,5^{2}\,7\cdot11 ,
\]
independently of the bracket monomial used.  The same construction at
$\delta=18$ gives $\Phi_{18}(\det_3)/K'_6=4725=3^{3}5^{2}7$ and
$\Phi_{18}(\per_3)/K'_6=-272=-2^{4}\cdot17$, whose quotient is the ratio
$-4725/272$ of \eqref{eq:value} and \eqref{eq:perm}.  Note also that
$K_a=K_{8-a}$ pairs indices of \emph{equal} parity, so it does not annihilate
the alternating sum: no symmetry of $U$ forces \eqref{eq:value24} to vanish.
This was recorded as a pre-registered prediction before the value was computed.
\end{remark}

\begin{remark}[Why the divisor does not decide it]\label{rem:notdivisor}
It is tempting to settle degree $24$ from Proposition \ref{prop:divisor} without
any computation.  On the orbit there is an invariant $\varphi$ of degree $6$
with $\varphi(g\cdot\det_3)=\det(g)^{2}$, well defined because $\det(h)=\pm1$
for $h\in H$, and $\Phi_{18}|_{G\cdot\det_3}=\Phi_{18}(\det_3)\varphi^{3}$; so
$\operatorname{div}(\varphi)=2P_1+3P_2$, and $\varphi^{4}$ has effective
divisor.  The inference ``effective divisor, hence regular'' is Serre's
criterion and needs $\cO$ normal, which it is not; applied to $\varphi$ itself
the same inference would give $6\in\mathcal E(\det_3)$, which is false.  What
the computation does prove is the positive part: $\varphi$ is regular on the
normalisation $\cO^{\nu}$, so $\bC[\cO^{\nu}]^{\SL_9}=\bC[\varphi]$ has degree
monoid exactly $6\bN$, and $\bC[\cO]^{\SL_9}$ is the semigroup ring of
$S\subseteq6\bN$ inside it.  The gaps of $\tfrac16S$ --- now known to be
$\{1,2\}$ or $\{1,2,5\}$ --- are precisely the degrees at which the
non-normality of $\cO$ is arithmetically visible on invariants, and the
conductor of that numerical semigroup is the conductor of $\bC[\cO]^{\SL_9}$ in
its normalisation.
\end{remark}

\begin{corollary}\label{cor:perm}
$\per_3\notin\cO$; that is, the border determinantal complexity of the
$3\times3$ permanent exceeds $3$.
\end{corollary}

\begin{proof}
By Theorem \ref{thm:ambient} the ambient space of degree-$18$ invariants is
one-dimensional, so $\Phi_{18}$ is defined on all of $\Sym^{3}\bC^{9}$ and its
value at any point is canonical once the scale is fixed; no choice of lift is
involved.  Evaluating at the permanent gives
\begin{equation}\label{eq:perm}
  \Phi_{18}(\per_3)\;=\;+50\,536\,120\,320\;=\;2^{20}\,3^{4}\,5\cdot7\cdot17
  \;\neq\;0 .
\end{equation}
By Proposition \ref{prop:divisor}, $\Phi_{18}$ vanishes identically on
$\partial\cO$, so $\per_3\notin\partial\cO$.  If $\per_3$ lay in $\cO$ it would
therefore lie in the dense orbit $G\cdot\det_3$, contradicting the theorem of
Marcus and Minc \cite{MM} that no element of $\GL_9$ carries $\det_3$ to
$\per_3$.  Hence $\per_3\notin\cO$.
\end{proof}

The conclusion of Corollary \ref{cor:perm} is of course classical and easy by
other means; what is new is the route.  The obstruction is a single integer,
computed once, and the argument uses no asymptotics, no representation-theoretic
multiplicity comparison, and no degeneration analysis --- exactly the shape of
separation that geometric complexity theory asks for --- and the shape in which
the strongest conditional lower bounds are stated \cite{LR} --- at the smallest
scale
where the shape can be exhibited at all.

\begin{remark}[Comparison]\label{rem:BI}
The determinant is the case where the general theory stops short.  For even
$n$, $e(\det_n)=n^{2}$ holds if and only if a signed count of admissible
$n$-tables is nonzero \cite[Prop.~3.28]{BI}, an Alon--Tarsi-type criterion
verified at $n=2$ and, by computer, at $n=4$; the analogous statement for the
product of variables is Kumar's \cite{KumarCMH,KL}.  For \emph{odd} $n$ the
criterion has no analogue --- the degree-$n^{2}$ invariant on the ambient space
vanishes identically for parity reasons --- so nothing in the literature
predicts a value, and $e(\det_n)$ was known for no odd $n$.  Theorem
\ref{thm:census} supplies the first one.  Whether $18=2n^{2}$ reflects a
pattern or is an accident of $n=3$ we do not know; the only other exact values
are $e(\det_2)=4$ and $e(\det_4)=16$, both equal to $n^{2}$, and three points
are not a sequence.  See Section \ref{sec:open}.
\end{remark}

\subsection*{The first conductor of a determinant}

With $\Phi_{18}$ as boundary semiinvariant, Definition \ref{def:deficit}
applies verbatim: the conductor ray attached to $(\lambda,\delta)$ is
$\big(\lambda+6k\mathbf 1_9,\ \delta+18k\big)$.

The two halves of the next statement have very different characters, and we
separate them because conflating them misrepresents what rests on what.

\begin{proposition}[The deficit]\label{prop:def222}
$\dfc\big((2,2,2),2\big)=1$.
\end{proposition}

\begin{proof}
Elementary in both factors.  For the closure multiplicity,
$\mult_\lambda\bC[\cO]_2\le\mult_\lambda\Sym^{2}\big(\Sym^{3}\bC^{9}\big)$, and
the classical plethysm $\Sym^{2}(\Sym^{3})=S_{(6)}\oplus S_{(4,2)}$ contains no
$S_{(2,2,2)}$ whatever; hence $\mult_\lambda\bC[\cO]_2=0$ and
$\dfc=m(\lambda)$.  For $m(\lambda)$: the identity component of $H$ acts through
$\bC^{9}=\bC^{3}\otimes\bC^{3}$, so an invariant in
$S_{(2,2,2)}(V\otimes W)$ must lie in $S_\mu V\boxtimes S_\nu W$ with both
factors one-dimensional, forcing $\mu=\nu=(2,2,2)$; the multiplicity is the
Kronecker coefficient $g\big((2,2,2),(2,2,2),(2,2,2)\big)=1$.  The residual
$\bZ/2$ generated by transposition acts on that line by a sign, which is $+1$.
Hence $m(\lambda)=1$.
\end{proof}

No orbit point, no evaluation functional, and no computation beyond a Kronecker
coefficient enters that proof.  The conductor is a different matter.

\begin{theorem}[The conductor]\label{thm:detcond}
$c\big((2,2,2),2\big)=1$, and the vanishing is \emph{ray-complete}:
$\dfc\big(\lambda+6k\mathbf 1_9,\ 2+18k\big)=0$ for every $k\ge1$.
\end{theorem}

Here the certification is genuinely needed, and here is what it certifies.  The
$k=1$ rung asks whether $\mult_\lambda\bC[\cO]_{20}$ attains $m(\lambda)$, and
the witness is the non-vanishing of one explicit function: the
$\lambda'=(8,8,8,6^{6})$ highest weight vector in degree $\delta=20$, whose
value is the $\mathrm{TOTAL}$ of Section \ref{sec:functional}.  It was evaluated
independently at two points of $\cO$, called $C$ and $R$, which Lemma
\ref{lem:displ} shows lie in distinct $H$-double cosets; both returned
$1\,152\,144\,000\ne0$.  A point-evaluation cannot certify a deficit, which is a
difference of dimensions; what it certifies is that a specific function is not
identically zero, hence that a multiplicity is at least one.  Ray-completeness
for $k\ge2$ then follows from the stabilisation argument of Section
\ref{sec:setup} together with the $k=1$ result.

This is the first conductor of a determinant to be computed exactly, and it is
the smallest possible nonzero answer: one multiplication by the fundamental
invariant clears the gap, and it stays cleared for ever after.  The scales are
worth noting --- the deficit is detected at degree $2$, far below the degree
$18$ at which the invariant that repairs it first exists.

\section{The evaluation functional}\label{sec:functional}

The computations of Section \ref{sec:det} all pass through one object: a
$\bZ$-valued functional on the orbit, evaluated at a point $u$ presented in
\emph{slab normal form}.  Writing a point of $\bC^{3}\otimes\bC^{3}\otimes\bC^{3}$
as a triple of $3\times3$ slabs and normalising the first slab to the identity,
a normal form is a pair $N=(A,B)$ of $3\times3$ matrices, and the residual
symmetry is simultaneous conjugation $(A,B)\mapsto(gAg^{-1},gBg^{-1})$ together
with the finite symmetries of the slab frame.

\begin{lemma}[Double cosets]\label{lem:displ}
Write points of the orbit as $g\cdot\det_3$ with $g\in G$.  The value of the
functional at $g\cdot\det_3$ depends only on the double coset $HgH\in
H\backslash G/H$.
Moreover the \emph{intersection algebra}
\[
  \mathfrak s(g)\;=\;\mathfrak h\cap\Ad(g)\,\mathfrak h
  \;\subseteq\;\mathfrak{gl}_9,
  \qquad \mathfrak h=\operatorname{Lie}H,
\]
is carried to $\Ad(h)\mathfrak s(g)$ under $g\mapsto hgh'$ for $h,h'\in H$, so
its isomorphism type is an invariant of the double coset.  Consequently two points whose intersection algebras are non-isomorphic
lie in different double cosets, and a value certified at one is an independent
certification when reproduced at the other.
\end{lemma}

We use this in exactly one place, and it is the place that matters.  At the two
certification points of Theorem \ref{thm:detcond},
\[
  \dim\mathfrak s(g_C)=\dim\mathfrak s(g_R)=4,
\]
but $\mathfrak s(g_C)$ is non-abelian while $\mathfrak s(g_R)$ is abelian.  The
two points are therefore genuinely inequivalent, and the agreement
\[
  \mathrm{TOTAL}(g_C)\;=\;\mathrm{TOTAL}(g_R)\;=\;1\,152\,144\,000
\]
is a check with content rather than a repetition of one computation in two
coordinate systems.  This is the strongest form of internal verification
available to us: no symmetry of the setup carries one point to the other, so
the shared value is a statement about the functional and not about the frame.
Theorem \ref{thm:totals} now explains the agreement --- the two points have the
same value of the gauge --- but it does not weaken the check, since the
theorem's normalisation is fixed at $C$ and the value at $R$ was measured
before any of it existed.

\begin{proposition}[Taxonomy of vanishing]\label{prop:taxonomy}
Every vanishing subvalue in the census of Section \ref{sec:det} is accounted
for by exactly one of three mechanisms, and the three are mutually exclusive on
the cases that occur:
\begin{enumerate}
\item[(i)] \emph{mass infeasibility} --- no admissible assignment of level
masses exists and the sum is empty.  This tier is decided by the
\emph{displacement-cancellation lemma}: content is feasible if and only if
there exist $2+2$ conversions with $\sum(e_{\mathrm{tgt}}-e_{\mathrm{src}})=0$,
which follows from the forced balanced distribution $t=(2,2)$ together with an
integral Birkhoff decomposition for $3\times3$ doubly stochastic matrices.  It
is a proof, not a test, and it derives the observed $15$-of-$81$ sweep
structurally;
\item[(ii)] \emph{wedge structure} --- assignments exist but every one of them
carries a repeated factor in an antisymmetrised slot, so each term vanishes
individually;
\item[(iii)] \emph{sign cancellation} --- terms exist and are individually
nonzero, and the sum vanishes.
\end{enumerate}
Only (iii) requires the computation; (i) and (ii) are decided by inspection of
the mask, and together they account for the great majority of the census.
\end{proposition}

The taxonomy is what makes the census tractable: it converts most of the
vanishing into combinatorics that can be certified without arithmetic, and it
isolates the cases where cancellation is doing real work.  It also has a
diagnostic use --- a subvalue that vanishes for reason (iii) at one evaluation
point and for reason (i) at another is a signal that the two points are being
compared incorrectly, and this caught two errors during development.

\subsection*{The gauge}

The totals at distinct normal forms are not equal, but they are constrained,
and the constraint turns out to be representation-theoretic.  Write
$\mathcal{B}$ for the space of bidegree-$(2,2)$ simultaneous-conjugation
invariants of a pair of $3\times3$ matrices.  Inside $\mathcal{B}$ put
\[
\begin{aligned}
  u_1 &= \tr(A^{2})\tr(B^{2})-\tr(AB)^{2}, \\
  u_2 &= \tr(A^{2}B^{2})-\tr\big((AB)^{2}\big), \\
  D   &= \big(\tr A\,\tr B-\tr(AB)\big)^{2}
        -\big((\tr A)^{2}-\tr A^{2}\big)\big((\tr B)^{2}-\tr B^{2}\big),
\end{aligned}
\]
and set $v:=u_1-2u_2$ and
\begin{equation}\label{eq:psi}
  \Psi\;:=\;2v-D\;=\;2u_1-4u_2-D .
\end{equation}
The third generator is where the obvious guess fails:
$\big(\tr A\,\tr B-\tr(AB)\big)^{2}$ by itself does not vanish on proportional
pencils, and the discriminant completion $D$ is the correct element.  On the
slot coordinates $(u_1,u_2,D)$ the sampled points read $C=(-1,-1,1)$,
$R=(1,0,1)$, $T_4=(3,0,5)$, the rank-one point $P=(0,0,0)$, and the entire
compression family $(2,0,4)$; so $\Psi=1$ at $C$, $R$ and $T_4$, and $\Psi=0$
at $P$ and on all of compression.

A warning about names, since they were assigned at different times and do not
follow one convention.  The subscript of $T_4$ counts \emph{transvections} ---
it is the four-transvection point $x_3\!+\!=\!x_0$, $x_4\!+\!=\!x_1$,
$x_7\!+\!=\!x_1$, $x_8\!+\!=\!x_2$ --- and $\Psi(T_4)=1$.  The subscripts of
$X_4$ and $X_{-3}$ record their \emph{gauge values}, $\Psi=4$ and $\Psi=-3$.
$T_4$ and $X_4$ are different points with different gauge values; we keep the
names because the repository and the value records use them.

\begin{proposition}[Uniqueness]\label{prop:unique}
$\dim\mathcal{B}=9$: the ten trace words of bidegree $(2,2)$ satisfy exactly
one relation, the polarised Cayley--Hamilton identity in rank $3$.  Inside
$\mathcal B$, slab-equivariance with character $\det^{2}$ --- that is
$\mathcal{D}_{A}f+2\tr(A)\,f=0$, and likewise in $B$ --- cuts out a
\emph{one}-dimensional space, spanned by $\Psi$.  In particular the
coefficient $-4$ on $u_2$ in \eqref{eq:psi} is forced, not fitted.
\end{proposition}

Two properties follow from the shape alone, with no computation.  Any
$\det^{2}$-equivariant covariant vanishes identically on proportional pencils,
since $(A,cA)$ is a degenerate $\GL_2$-transform of $(A,0)$; this accounts for
the exact zeros at $P$.  And $\Psi$ is shift-invariant, which is what makes the
normal form well posed.

\begin{theorem}[The gauge is the Aronhold invariant]\label{thm:gauge}
Let $I_6$ be the fundamental degree-$6$ invariant of
$\bC^{3}\otimes\bC^{3}\otimes\bC^{3}$ under
$\SL_3\times\SL_3\times\SL_3$ --- the lowest of the three basic invariants, of
degrees $6$, $9$ and $12$, in Vinberg's classification of this
$\theta$-representation \cite{Nur,BHO}.  Then
\begin{equation}\label{eq:i6}
  I_6(I,A,B)\;=\;-6\,\Psi(A,B),
\end{equation}
an identity of polynomials in all $18$ pencil indeterminates.  Equivalently
$\Psi=-I_6^{\mathrm{prim}}(I,A,B)$ with coefficient $1$, where
$I_6^{\mathrm{prim}}$ is the primitive integral normalisation.
\end{theorem}

This is not the statement that $\Psi$ is one component of the restriction: the
restriction has no other component.  Each $\varepsilon$ in the first slot needs
three distinct slab indices, so every three-block uses $I$, $A$ and $B$ once
each and the slab multidegree is forced to $(2,2,2)$.  The content of $I_6$ is
$6$, which is the factor in \eqref{eq:i6}.  As an independent check on the
identification, our $I_6$ has exactly $1152$ monomials on the generic
$3\times3\times3$ tensor, matching the Aronhold invariant as tabulated by
Bremner and Hu \cite{BHO}.

\begin{theorem}[Totals law]\label{thm:totals}
For every balanced direction $N=e_1\otimes A+e_2\otimes B$,
\[
  \mathrm{TOTAL}(N)\;=\;\Psi(N)\cdot 1\,152\,144\,000 .
\]
\end{theorem}

The measured totals are $1\,152\,144\,000$ at $C$ and at $R$, where $\Psi=1$;
$4\,608\,576\,000$ at $X_4$, where $\Psi=4$; and $0$ at the rank-one point $P$
and at each of four compression points, where $\Psi=0$.  The law was stated as
an empirical identity in an earlier version of this paper, tested at
$\Psi=1,4,0$ and at no other value; the proof below closes it.  The one value it
predicts that has never been measured is at $X_{-3}$, where $\Psi=-3$, and we
record the prediction $\mathrm{TOTAL}(X_{-3})=-3\,456\,432\,000$ as a standing
test of the two inputs the proof consumes.

For orientation: the functional is a $\GL_9$-Borel highest weight vector in
$\Sym^{20}\Sym^{3}\bC^{9}$ --- degree $20$, not $18$, because it lives at the
$k=1$ rung $\delta=2+18$ of the conductor ray of Theorem \ref{thm:detcond}, not
in the degree of the invariant.

\subsection*{Proof of the totals law}

Three facts go in, and the third is Proposition \ref{prop:unique}.  The first
two are recorded here.

\begin{lemma}[Inputs]\label{lem:inputs}
\emph{(a)}  Content forces exactly four substituted legs, distributed as
$t=(2,2)$, so $\mathrm{TOTAL}$ is a polynomial of bidegree $(2,2)$ on the
$18$-dimensional space of balanced directions
$N=e_1\otimes A+e_2\otimes B$.
\emph{(b)}  Let $Q\subset\GL_9$ be the stabiliser of the highest-weight line of
$h$, a parabolic with Levi $\GL_3\times\GL_6$ and $Q=\mathrm{Stab}(W)$ for the
$6$-dimensional coordinate subspace $W$, and let $\chi$ be its character on that
line, $\chi(q)=\det(q|_{V/W})^{8}\det(q|_{W})^{6}$.  If $u_{N'}=q\,u_N\,s$ with
$q\in Q$ and $s\in H$ then $\mathrm{TOTAL}(N')=\chi(q)\,\mathrm{TOTAL}(N)$.
Moreover $\Gamma_N:=u_N^{-1}(W)$ satisfies $\Gamma_{N'}=s^{-1}\Gamma_N$, so the
double coset of $u_N$ is exactly the $H$-orbit of the net $\Gamma_N^{\perp}$
\emph{as a subspace} of $M_3$; in the normal form of this section
$\Gamma_N^{\perp}=\{[v;Av;Bv]:v\in\bC^{3}\}$, whose tensor slabs are $(I,A,B)$.
\end{lemma}

What was missing was the identification of $\chi$, and it has two halves.

\begin{lemma}[The character is a square]\label{lem:sq}
Write $s=t^{-1}$, so that $q=u_{N'}\,t\,u_N^{-1}$ and $\Gamma_{N'}=t\,\Gamma_N$.
Then
\[
  \chi(q)=\det(q|_{V/W})^{8}\det(q|_{W})^{6}
        =\det(q)^{6}\,\det(q|_{V/W})^{2}
        =\det(q|_{V/W})^{2}.
\]
\end{lemma}

\begin{proof}
The first equality is $8=6+2$ and $\det q=\det(q|_{V/W})\det(q|_{W})$.  For the
second, $u_N$ and $u_{N'}$ are unipotent, so $\det q=\det t$; and $\det t=\pm1$
on $H$, since $X\mapsto aXb$ has determinant $(\det a\det b)^{3}=1$ and the
transpose has determinant $(-1)^{3}=-1$ on $M_3$.  Hence $\det(q)^{6}=1$ on
both cosets of $H$.
\end{proof}

The transpose coset therefore cannot contribute to $\chi$ --- the same parity
that closes the ray in Theorem \ref{thm:detcond}, where $\det{}_9$ is $-1$ on
the transpose coset and the exponent is again even.

\begin{lemma}[Slot dictionary]\label{lem:slot}
$(V/\Gamma_N)^{*}=\Gamma_N^{\perp}$, and the basis of $\Gamma_N^{\perp}$
transported from $W^{\perp}$ by $u_N$ is exactly the canonical parameter $v$ of
the normal form $\Gamma_N^{\perp}=\{[v;Av;Bv]\}$.  Consequently
$\det(q|_{V/W})=\det(G)^{-1}$, where $G$ is the change of parameter induced by
$t$ --- that is, the basis change in the \emph{third tensor slot}, the one on
which $H$ does not act.  Explicitly, with $\alpha=a^{-T}$ and $\beta=b^{-T}$,
\[
  G^{T}=\beta^{T}(\alpha_{00}I+\alpha_{01}A+\alpha_{02}B),\qquad
  A'=\bigl(\beta^{T}(\alpha_{10}I+\alpha_{11}A+\alpha_{12}B)\bigr)(G^{T})^{-1},
\]
and likewise for $B'$.  Hence $\chi=\det(G)^{-2}$.
\end{lemma}

The two lemmas together are the identification $\chi\leftrightarrow\det^{2}$:
the character is $\det^{2}$ in the net's parameter slot.  Since $I_6$ is a
semi-invariant of character $\det^{2}$ in each slot and the two $M_3$-side
factors contribute $(\det a\det b)^{-2}=1$, the gauge obeys the same law,
\[
  \Psi(A',B')=\det(G)^{-2}\Psi(A,B)=\chi\cdot\Psi(A,B),
\]
on both cosets of $H$ --- on the transpose coset because the transpose acts on
the net by $Y\mapsto Y^{T}$, which swaps the first and third tensor slots, and
$S_3$ acts on the one-dimensional degree-$6$ invariant space through the
\emph{trivial} character.  (Had it been the sign character, $\Psi$ would be
anti-invariant where $\mathrm{TOTAL}$ is invariant and the law would be false.
It is not: $\Psi$ does not vanish on the nine-parameter family of tensors fixed
by that swap.)

\begin{proof}[Proof of Theorem \ref{thm:totals}]
Specialise $t$.  Taking $\alpha=I$ and $\beta\in\SL_3$ gives $G=\beta^{T}$ with
$\det G=1$, so $\chi=1$ and $\mathrm{TOTAL}$ is a simultaneous-conjugation
invariant.  Taking $\beta=I$ and $\alpha$ the unipotent sending
$\mathrm{slab}_0\mapsto\mathrm{slab}_0+t\,\mathrm{slab}_1$ gives
$(A,B)\mapsto\bigl(A(I+tA)^{-1},B(I+tA)^{-1}\bigr)$ with
$\chi=\det(I+tA)^{-2}$, whose infinitesimal form is
$\mathcal D_A\mathrm{TOTAL}+2\tr(A)\mathrm{TOTAL}=0$; the $\mathrm{slab}_2$
mirror gives the same in $B$.  Both elements lie in $H$, since $\alpha$ and
$\beta$ are unimodular.  So $\mathrm{TOTAL}$ lies in $\mathcal B$ and is
slab-equivariant with character $\det^{2}$; by Proposition \ref{prop:unique}
that space is one-dimensional and spanned by $\Psi$, whence
$\mathrm{TOTAL}=c\,\Psi$.  Evaluating at $C$, where $\Psi=1$ and the measured
total is $1\,152\,144\,000$, gives $c$.
\end{proof}

Two remarks on what the proof explains.  First, $C$ and $R$ lie in
\emph{different} double cosets, so no symmetry of the setup carries one to the
other; their equal totals looked like a coincidence and are now forced, because
the group is too small to connect the points but the space of functions it
constrains is one-dimensional.  Second, the vanishing of the totals on the
compression family is the classical degeneracy of the Aronhold invariant rather
than a feature of the evaluation.

\begin{remark}[Why $\det^{2}$, and what happens at other weights]\label{rem:pq}
Nothing above is special to $\Psi$.  For a deficit weight
$\lambda'=(p,p,p,q^{6})$ the same computation gives
$\chi=\det(t)^{q}\det(q|_{V/W})^{p-q}$, so the character in the parameter slot
is $\det^{p-q}$ and the transpose coset is harmless precisely when $q$ is even.
Here $p-q=8-6=2$, which is why the governing object is the degree-$6$ Aronhold
invariant; at a weight with $p-q=3$ it would be the degree-$9$ Vinberg
generator, and at $p-q=4$ the degree-$12$ one.  When $q$ is odd, $\chi$ acquires
$-1$ on the transpose coset while any $\det^{\mathrm{even}}$ semi-invariant does
not, which forces $\mathrm{TOTAL}$ to vanish on every direction whose net is
$H$-equivalent to its own transpose.  We have tested none of this.
\end{remark}

\begin{remark}[An arithmetic signature]\label{rem:arith}
Every value the functional has returned, at every point and every $\sigma$, is
an integer multiple of $151\,200=2^{5}3^{3}5^{2}7$.  The reader can check this
on the one table we give in full.  At $X_4$ the eight orbit representatives
carry weights $(4,4,8,4,8,4,2,2)$, summing to the $36$ subproblems, and
cofactors
\[
  -2038,\ 4338,\ 3843,\ -258,\ 3567,\ -4372,\ -5188,\ -4552 ,
\]
whose weighted sum is $30\,480=\mathrm{TOTAL}(X_4)/151\,200$.  At $C$, $R$, $Q$
and $T_4$ the corresponding leading cofactor is $719$, and at $X_{-3}$ it is
$5907$.  The prime $719$ organises the whole table and is what makes the
refutation in Remark \ref{rem:fails}(a) sharp: it divides no other cofactor, so
a proportionality $V_\sigma=W(\sigma)\Psi$ would force a non-integer where the
engine returns an integer.

The common factor is not imposed anywhere in the computation, and we have no
proof of it.  Two further coincidences we flag without explaining:
$\mathrm{TOTAL}(C)=2^{7}3^{4}5^{3}\cdot7\cdot127$ and
$\mathrm{TOTAL}(X_4)=2^{9}3^{4}5^{3}\cdot7\cdot127$, so the prime $127$ sits in
the totals constant and nowhere else in the paper; and
$\Phi_{18}(\det_3)=-2^{16}3^{7}5^{3}7^{2}$ is $151\,200$ times
$-5\,806\,080$, but it is a value in degree $18$ and belongs to a different
computation from the degree-$20$ table above, so we draw no inference from the
shared factor.
\end{remark}

\begin{remark}[What fails]\label{rem:fails}
Three hypotheses about the functional were pre-registered and then refuted by
their own tests.  We record them because each was natural, each shaped the
search, and each was killed by a single number.
\emph{(a)} The individual matrix elements $V^{h}_{\sigma}$ do \emph{not}
inherit $\GL_2$-covariance: the engine returned $f1X4_{00}=-308\,145\,600$
against a pre-registered $+434\,851\,200$, gate-confirmed by an independent
run.  The highest weight vector is a $\GL_9$-Borel object and the choice of
scheme and of $\sigma$ breaks the covariance.
\emph{(b)} The per-$\sigma$ values are not simultaneous-conjugation invariants
either.  A parameter-free fit over $\mathcal B$ reached rank $9$, the maximum
on the feasible cone, and predicted $f1Y4_{00}=+69\,854\,400$; the engine
returned $0$.  The violated relation spreads over nearly every point, so it is
not an outlier.
\emph{(c)} Both plane-cubic routes to the gauge are dead --- there are no
degree-$2$ invariants of ternary cubics, and three triangles realise
$\Psi=0,1,4$ --- and the pencil cubic is insufficient, since $\Psi$ requires
the length-four word $u_2$.
None of this touches the results of Section \ref{sec:det}: the conductor, the
value at $R$, and the ray closure use no covariance assumption anywhere.  What
the refutations do is confine the phenomenon to the totals level, which is
where Conjecture \ref{thm:totals} states it.
\end{remark}

\section{Verification protocol}\label{sec:method}

Several of the results above are computational, and we regard the protocol as
part of their content rather than as an appendix to it.

\smallskip
\noindent\textbf{Exact arithmetic throughout.}  No floating point occurs
anywhere in the pipeline.  Multiplicities are computed by an exact streamed
level algorithm over $\bZ$: level states are packed into $54$-bit masks held in
machine words, the working set is an open-addressing table, and spilled runs
are written to disk in a delta-varint code and merged.  When the working set
exceeds the memory budget the level is re-run in shards, with the shard count
doubling until it fits; the answer is identical, only the schedule changes.
Checkpoints are written atomically at every spill, so an interrupted run
resumes at the last spill rather than at the last level.

\smallskip
\noindent\textbf{Independent routes.}  The census values of Theorem
\ref{thm:census} and the totals of Theorem \ref{thm:gauge} were each produced by
two routes that share no code below the interface: the direct level algorithm,
and a compressed route that quotients by the symmetry harvest first.  The
compression is not a heuristic --- that the compressed sum reproduces the full
sum is proved, by exhibiting the sign character on each orbit of subproblems
and checking that the harvest is a group action with the claimed orbit sizes.
The harvest reduces $36$ subproblems to $18$ and then to $11$ inequivalent
runs; the reduction is exact and the two routes agree on every value reported
here.

\smallskip
\noindent\textbf{Anchors.}  The pipeline is pinned at both ends by cases with
independently known answers.  At the small end, the same code path run on
$\bC^{2}\otimes\bC^{2}\otimes\bC^{2}$ reproduces Cayley's hyperdeterminant
\cite{GKZ} on the standard test points: $0$ at the two rank-one points, $0$ at
the $W$-state (rank $3$, on the tangent locus), $1$ at the GHZ point, and
$9\,572\,836=3094^{2}=\big(\det(u,u')\det(v,v')\det(w,w')\big)^{2}$ at a
generic rank-two point --- exactly, not to within rounding.  This is an
independent check on the invariant-theoretic layer, not merely on the
arithmetic.  At the classical end, the closed forms of Theorems
\ref{thm:A} and \ref{thm:B} are reproduced by the general machinery.

\smallskip
\noindent\textbf{Canonical regression.}  A fixed regression suite is re-run
before every reported value.  Its published values are: the quadratic anchors
$\mathrm{quad}=24$, $\mathrm{quad}_0=0$, and $\mathrm{quad}_q$ raw $6\times4=24$;
the $\det_3$ ladder $L_2=29/29/29$ through
$L_6=1818118/2336283/2686868$; and, at the certification point,
$L_7=54685987/100774838/141001840$ and
$L_8=128027708/422952740/603408404$.  Any deviation in any entry aborts the
run.  These numbers are stated here so that a reimplementation can be checked
against them without reproducing the whole computation.

\smallskip
\noindent\textbf{Pre-registration.}  Predictions were committed to version
control before the corresponding values existed, and the commit history is
part of the record: for each of the results above one can check what was
predicted, when, and whether it held.  Two of the hypotheses so registered were
refuted (Remark \ref{rem:fails}), and we consider their presence in the record
to be evidence about the method rather than an embarrassment to it.  A claim
that was formed after seeing the value is marked as such wherever it appears.

\section{Open problems}\label{sec:open}

\begin{question}[Transport in general]\label{q:transport}
Theorem \ref{thm:B} computes the conductor from boundary data alone,
$c(\lambda)=\lfloor\mu_{\max}(\lambda)/|w_N|\rfloor$, on all $254$
deficit-positive weights in range.  Does this shape hold for every orbit
closure with a smooth dense boundary orbit and a one-dimensional normal weight?
The upper bound is the block-order estimate and is general; what is missing is
attainment, and Remark \ref{rem:orphan} shows attainment can fail exactly on
empty support.
\end{question}

\begin{question}[The totals law at other weights]\label{q:gauge}
Theorem \ref{thm:totals} is now proved, and Remark \ref{rem:pq} says why the
character is $\det^{2}$: for a deficit weight $\lambda'=(p,p,p,q^{6})$ one gets
$\chi=\det(t)^{q}\det(q|_{V/W})^{p-q}$, so the governing object should be the
$\SL_3^{3}$-invariant of degree $3(p-q)$ --- Aronhold's $I_6$ here, and the
degree-$9$ and degree-$12$ Vinberg generators when $p-q=3$ and $4$.  Is that so?
The question has two parts, and only the second needs a computation.  Does the
counting lemma still force a balanced bidegree at those weights, so that the
uniqueness step has a one-dimensional space to land in?  And when $q$ is odd,
does $\mathrm{TOTAL}$ really vanish on every direction whose net is
$H$-equivalent to its own transpose, as the parity in Lemma \ref{lem:sq}
demands?  A single evaluation at such a direction would settle the second.
\end{question}

\begin{question}[Degree $30$, and the shape of the semigroup]\label{q:deg30}
Corollary \ref{cor:semigroup} leaves exactly one degree open: is
$30\in\mathcal E(\det_3)$?  Equivalently, is $S=\langle18,24\rangle$ or
$S=\langle18,24,30\rangle$?  Unlike degree $24$ this is not decided by one
number: the ambient space $\bC[\Sym^3\bC^9]^{\SL_9}_{30}$ is
four-dimensional, so nonvanishing is still witnessed by a single evaluation but
vanishing would require all four.  Two remarks on what to expect.  Every
numerical semigroup with two generators is symmetric, so
$S=\langle18,24\rangle$ holds if and only if $\bC[\cO]^{\SL_9}$ is Gorenstein,
whereas $\langle3,4,5\rangle$ is not symmetric; we know no reason for
Gorensteinness here and offer this only as the shape of the dichotomy.  And by
Remark \ref{rem:notdivisor} the answer names the last gap of $\tfrac16S$, hence
the exact conductor of $\bC[\cO]^{\SL_9}$ in its normalisation $\bC[\varphi]$:
either $36$ or $24$.
\end{question}

\begin{question}[The next conductor row]\label{q:delta3}
Theorem \ref{thm:detcond} settles $\delta=2$.  The row $\delta=3$ requires
either a theoretical reduction or a substantially larger computation; we
believe the theory should come first, since the pattern of Section
\ref{sec:worldB} suggests the answer is again a floor and the interesting
content is the numerator.
\end{question}

\begin{question}[$n=4$ and beyond]\label{q:n4}
The three known values are $e(\det_2)=4$, $e(\det_3)=18$ and $e(\det_4)=16$.
The even ones sit at $n^{2}$ and the odd one does not, which suggests the
parity split that runs through \cite{BI} is the real division rather than any
growth pattern in $n$.  The first test is $n=5$: is $e(\det_5)>25$, and is it
again $2n^{2}$?  We expect the degree period to do as much work there as it did
here --- $5\equiv1\pmod4$ gives $a=1$ and $b=5$, a coarser constraint than we
enjoyed at $n=3$, so the computation would be correspondingly harder.  Closer
to hand: what is the first deficit-positive weight of $\det_4$, and what is its
conductor?  A negative answer to Question \ref{q:transport} would most likely
be visible there first.
\end{question}

\section*{Data availability}

All code, inputs, exact outputs, and the full commit history --- including the
pre-registration commits referred to in Section \ref{sec:method} --- are
available at
\begin{center}
  \url{https://github.com/swsethuraman/gct}
\end{center}
The regression values quoted in Section \ref{sec:method} are reproduced by the
scripts in that repository without any of the long runs.

\section*{Acknowledgements}

The work reported here was carried out jointly with Claude (Anthropic).
Responsibility for the correctness of every statement in this paper rests with
the named author.

\begin{thebibliography}{99}

\bibitem{AFPRW}
M.~Aprodu, G.~Farkas, {\c S}.~Papadima, C.~Raicu, J.~Weyman,
\emph{Koszul modules and Green's conjecture},
Invent. Math. \textbf{218} (2019), 657--720.

\bibitem{BHO}
M.~R.~Bremner, J.~Hu, L.~Oeding,
\emph{The $3\times3\times3$ hyperdeterminant as a polynomial in the fundamental
invariants for $\SL_3(\bC)\times\SL_3(\bC)\times\SL_3(\bC)$},
Math. Comput. Sci. \textbf{8} (2014), 147--156; arXiv:1310.3257.

\bibitem{BI}
P.~B\"urgisser, C.~Ikenmeyer,
\emph{Fundamental invariants of orbit closures},
J. Algebra \textbf{477} (2017), 390--434.

\bibitem{BIP}
P.~B\"urgisser, C.~Ikenmeyer, G.~Panova,
\emph{No occurrence obstructions in geometric complexity theory},
J. Amer. Math. Soc. \textbf{32} (2019), 163--193.

\bibitem{GKZ}
I.~M.~Gelfand, M.~M.~Kapranov, A.~V.~Zelevinsky,
\emph{Discriminants, Resultants, and Multidimensional Determinants},
Birkh\"auser, Boston, 1994.

\bibitem{HL}
J.~H\"uttenhain, P.~Lairez,
\emph{The boundary of the orbit of the $3\times3$ determinant polynomial},
C. R. Math. Acad. Sci. Paris \textbf{354} (2016), 931--935; arXiv:1512.02437.

\bibitem{DM}
H.~Derksen, V.~Makam,
\emph{Polynomial degree bounds for matrix semi-invariants},
Adv. Math. \textbf{310} (2017), 44--63.

\bibitem{Hue}
J.~H\"uttenhain,
\emph{Geometric Complexity Theory and Orbit Closures of Homogeneous Forms},
Ph.D. thesis, Technische Universit\"at Berlin, 2017.

\bibitem{IK}
C.~Ikenmeyer, U.~Kandasamy,
\emph{Implementing geometric complexity theory: on the separation of orbit
closures via symmetries},
Proc. 52nd ACM Symposium on Theory of Computing (STOC 2020), 713--726.

\bibitem{KumarCMH}
S.~Kumar,
\emph{Geometry of orbits of permanents and determinants},
Comment. Math. Helv. \textbf{88} (2013), 759--788.

\bibitem{KL}
S.~Kumar, J.~M.~Landsberg,
\emph{Connections between conjectures of Alon--Tarsi, Hadamard--Howe, and
integrals over the special unitary group},
Discrete Math. \textbf{338} (2015), 1232--1238.

\bibitem{Kraft}
H.~Kraft,
\emph{Geometrische Methoden in der Invariantentheorie},
Aspekte der Mathematik D1, Vieweg, Braunschweig, 1984.

\bibitem{Lands}
J.~M.~Landsberg,
\emph{Geometry and Complexity Theory},
Cambridge Studies in Advanced Mathematics 169, Cambridge Univ. Press, 2017.

\bibitem{LR}
J.~M.~Landsberg, N.~Ressayre,
\emph{Permanent v. determinant: an exponential lower bound assuming symmetry},
Proc. 2016 ACM Conf. on Innovations in Theoretical Computer Science (ITCS),
29--35.

\bibitem{MM}
M.~Marcus, H.~Minc,
\emph{On the relation between the determinant and the permanent},
Illinois J. Math. \textbf{5} (1961), 376--381.

\bibitem{MS}
K.~D.~Mulmuley, M.~Sohoni,
\emph{Geometric complexity theory I: an approach to the P vs. NP and related
problems},
SIAM J. Comput. \textbf{31} (2001), 496--526.

\bibitem{Weyl}
H.~Weyl,
\emph{The Classical Groups: Their Invariants and Representations},
Princeton Univ. Press, 1939.

\bibitem{Nur}
A.~G.~Nurmiev,
\emph{Orbits and invariants of third-order matrices},
Sb. Math. \textbf{191} (2000), 717--724.

\end{thebibliography}

\end{document}
